\documentclass[11pt,aspectratio=169]{beamer}
\usepackage[T1]{fontenc}
\usepackage[utf8]{inputenc}
\usepackage[english]{babel}
\usepackage{mathpazo}
\usefonttheme{serif}
\usepackage{graphicx,microtype}
\usepackage{tikz}

\definecolor{blueI}{HTML}{2A78D6}
\definecolor{ink}{HTML}{0B0B0B}
\definecolor{steel}{RGB}{74,96,121}
\definecolor{steellight}{RGB}{192,206,220}

% Theme: steel-blue title bar, white pages, as in the other decks.
\setbeamercolor{structure}{fg=steel}
\setbeamercolor{normal text}{fg=ink,bg=white}
\setbeamercolor{frametitle}{fg=white,bg=steel}
\setbeamercolor{alerted text}{fg=blueI}
\setbeamercolor{panelhead}{fg=steel,bg=steellight!30}
\setbeamersize{text margin left=7mm,text margin right=7mm}
\setbeamertemplate{navigation symbols}{}
\setbeamertemplate{itemize item}{\textcolor{steel}{\small$\bullet$}}
\setbeamertemplate{itemize subitem}{\textcolor{steel}{\textendash}}
\setbeamertemplate{itemize/enumerate body begin}{\setlength{\itemsep}{0.3em}}
\setbeamertemplate{frametitle}{%
  \nointerlineskip
  \begin{beamercolorbox}[wd=\paperwidth,ht=9mm,dp=3mm,center]{frametitle}
    \large\insertframetitle\strut
  \end{beamercolorbox}}
\setbeamertemplate{footline}{}

% The header bar of one column.
\newcommand{\panelhead}[1]{%
  \begin{beamercolorbox}[wd=\linewidth,sep=5pt]{panelhead}\textbf{#1}\end{beamercolorbox}\par\vspace{2mm}}

\hypersetup{pdftitle={Running the viewer tool on different computers}}

\begin{document}

\begin{frame}[plain]
\begin{tikzpicture}[remember picture,overlay]
 \fill[steel](current page.south west)rectangle(current page.north east);
 \fill[steellight](current page.south west)rectangle([yshift=20mm]current page.south east);
 \node[text=white,align=center,text width=.9\paperwidth]at([yshift=8mm]current page.center)
   {{\LARGE\bfseries The viewer tool on other computers}\\[4mm]};
 \node[text=steel,align=center]at([yshift=10mm]current page.south)
   {October 2026};
\end{tikzpicture}
\end{frame}

\begin{frame}{Good news: The tool should now work on any computer}
\begin{columns}[T,onlytextwidth]
\column{.52\textwidth}
\panelhead{First time: two steps}
\begin{enumerate}
  \item install \texttt{uv} (and \texttt{git}) if missing: one command
  \item launch the viewer with the command below, and choose the data folder once
\end{enumerate}
\begin{itemize}
  \item no Python and no library to install by hand
\end{itemize}
\column{.44\textwidth}
\panelhead{Afterwards: one command}
\begin{itemize}
  \item the same launch command
  \item the data folder is remembered
  \item always the latest code from github.
\end{itemize}
\end{columns}
\vfill
\begin{beamercolorbox}[wd=\linewidth,sep=5pt,center]{panelhead}
\ttfamily\small uvx --from git+https://github.com/matteodisante/xc-viewer xc-viewer
\end{beamercolorbox}
\end{frame}

\begin{frame}{Bad news: running the viewer on the CFM laptop}
\begin{columns}[T,onlytextwidth]
\column{.36\textwidth}
\panelhead{1. Personal laptops}
\begin{itemize}
  \item install \texttt{uv}, one command
  \item simple and fast
  \item the data must sit on the laptop: about 120~GB today
\end{itemize}
\column{.60\textwidth}
\panelhead{2.CFM laptops}
\begin{itemize}
  \item \textbf{implemented}: one command fetches the code from github and installs the packages. Always up to date, but not compliant with the policy ?
  \item \textbf{compliant route?}: an environment (Python, libraries) requested for each laptop?
  \item even then, the code probably cannot come from the aforementioned github repo.
  \item a new feature needs new libraries: new authorizations, possibly refused. Development is limited.
\end{itemize}
\end{columns}
\vfill
\begin{beamercolorbox}[wd=\linewidth,sep=5pt]{panelhead}
Option 2 is operational, but I do not think it can be used on the CFM laptops.
\end{beamercolorbox}
\end{frame}

\begin{frame}{Developing the viewer on CFM laptops}
\begin{columns}[T,onlytextwidth]
\column{.48\textwidth}
\panelhead{Today: my personal Mac}
\begin{itemize}
  \item personal Claude Code, with the code of the viewer and of every analysis
  \item it knows the whole situation and reuses the functions/context of the analyses
  \item a new feature is simpler, faster, more effective
\end{itemize}
\column{.48\textwidth}
\panelhead{CFM laptops}
\begin{itemize}
  \item Claude without internet access
  \item no context of the other repositories, where the analyses are developed
  \item the personal github account cannot be used.
\end{itemize}
\end{columns}
\vfill
\begin{beamercolorbox}[wd=\linewidth,sep=5pt]{panelhead}
Without the context of the analyses, every new feature is harder to build.
\end{beamercolorbox}
\end{frame}

\end{document}